\documentclass[11pt]{article}
\usepackage[T1]{fontenc}
\usepackage[utf8]{inputenc}
\usepackage{lmodern,amsmath,amssymb,amsthm,mathtools}
\usepackage[margin=1in]{geometry}
\usepackage{microtype}
\usepackage[colorlinks=true,linkcolor=blue,citecolor=blue,urlcolor=blue]{hyperref}
\hypersetup{pdftitle={A uniform Kolmogorov bound for Jack measures},pdfauthor={Kaiwen Jin}}
\newtheorem{theorem}{Theorem}[section]
\newtheorem{lemma}[theorem]{Lemma}
\newcommand{\E}{\mathbb E}
\newcommand{\Prob}{\mathbb P}
\newcommand{\dK}{d_{\mathrm K}}
\newcommand{\fall}[2]{#1^{\downarrow #2}}
\title{A uniform Kolmogorov bound for Jack measures}
\author{Kaiwen Jin\\\small\href{mailto:yui.ui.siki@gmail.com}{yui.ui.siki@gmail.com}}
\date{October 13, 2026}
\begin{document}
\maketitle
\begin{abstract}
For the normalized transposition statistic under Jack measure on partitions
of $n$, we prove a Kolmogorov bound
$10\max\{n^{-1/2},\sqrt\alpha/n\}$ for every $n\ge3$ and $\alpha\ge1$.
This removes the logarithmic loss in the known bound, including the transition
$\alpha\asymp n$, and proves the uniform estimate conjectured by Fulman.
The proof retains cancellation in the conditional third moment of a Jack
growth exchangeable pair. Its expansion in three orthogonal Jack characters
gives the required parameter dependence.
\end{abstract}

\section{Introduction}

Jack measures deform Plancherel measure on integer partitions and provide
discrete analogues of Gaussian beta ensembles \cite{Fulman}.
For a partition $\lambda\vdash n$, put
\[
 c_\lambda=\prod_{u\in\lambda}(\alpha a(u)+\ell(u)+1),\qquad
 c'_\lambda=\prod_{u\in\lambda}(\alpha a(u)+\ell(u)+\alpha),
 \qquad
 \pi_{n,\alpha}(\lambda)=\frac{\alpha^n n!}{c_\lambda c'_\lambda}.
\]
Here $a(u)$ and $\ell(u)$ count the boxes to the right of and below $u$.
The parameter satisfies $\alpha>0$; at $\alpha=1$ the law is Plancherel
measure. Write $\lambda_i$ and $\lambda'_i$ for row and column lengths.
The statistic of interest is
\begin{equation}\label{eq:statistic}
 W_{n,\alpha}=
 \frac{\sum_i\left[\alpha\binom{\lambda_i}{2}-\binom{\lambda'_i}{2}\right]}
      {\sqrt{\alpha\binom n2}},\qquad \lambda\sim\pi_{n,\alpha}.
\end{equation}
It has mean zero and variance one \cite[Lemmas~3.1--3.2]{Fulman}.
At $\alpha=1$, it is the normalized
irreducible character ratio of a transposition. Its normal approximation
quantifies fluctuations of this character statistic as the partition size
and the deformation parameter vary.

Let $Z$ be standard normal, with distribution function $\Phi$, and write
\[
 \dK(V,Z)=\sup_{x\in\mathbb R}|\Prob(V\le x)-\Phi(x)|.
\]
Fulman conjectured a bound by a universal constant times
$\max\{n^{-1/2},\sqrt\alpha/n\}$ \cite[p.~4 of the author version]{Fulman}.
Chen, Rai\v c and Th\`anh proved
$8.2\max\{n^{-1/2},\sqrt\alpha\log n/n\}$
\cite[Proposition~4.1]{CRT}.
Their estimate has the desired order for $1\le\alpha\le n/\log^2n$;
they also obtained the desired order for $\alpha\ge n^{1+\varepsilon}$,
with a constant depending on the fixed $\varepsilon>0$.
Th\`anh and Tu's subsequent non-uniform estimate applies on the restricted
range $\log^2 n/n\le\alpha\le n/\log^2 n$ \cite[Theorem~8]{TT}.
The uniform Kolmogorov estimate across the intermediate range is the
remaining question. The following bound covers the whole parameter range.

\begin{theorem}\label{thm:main}
For every integer $n\ge3$ and every real $\alpha\ge1$,
\begin{equation}\label{eq:main}
 \dK(W_{n,\alpha},Z)
 \le 10\max\left\{n^{-1/2},\frac{\sqrt\alpha}{n}\right\}.
\end{equation}
\end{theorem}

The constant is independent of both parameters; no restriction is needed
when $\alpha$ is comparable to $n$. We do not optimize its numerical value.
The growth coupling of Kerov \cite{Kerov} supplies an exchangeable pair
whose increment need not have a small deterministic bound. We use the
normal-approximation inequality of Shao and Zhang \cite{SZ}, which depends
on the signed conditional third moment rather than the absolute third
moment of the increment. We compute this conditional moment as a linear
combination of three orthogonal Jack characters. Orthogonality then gives
its exact second moment and the uniform estimate.

\section{Jack growth and the cubic increment}

\subsection{Corner probabilities and content moments}

Fix $\alpha\ge1$ and write $d=\alpha-1$. The content of a box $(i,j)$,
with indices starting at one, is
$c(i,j)=\alpha(j-1)-(i-1)$. For any partition $\rho$, set
\[
 s_k(\rho)=\sum_{u\in\rho}c(u)^k.
\]
The contents sum to the numerator of \eqref{eq:statistic}, by summing $j-1$
along rows and $i-1$ along columns.
An outer corner has coordinate $x=c(u)$ for an addable box $u$.
An inner corner has coordinate $y=c(u)+d$ for a removable box $u$;
the shift arises because the inner coordinate uses the south-east corner
of the box. Define
\begin{equation}\label{eq:G}
 G_\rho(z)=\frac{\prod_{y\text{ inner}}(z-y)}{\prod_{x\text{ outer}}(z-x)}.
\end{equation}

Kerov's growth process adds a box to $\rho$ with content $X$, with
probabilities equal to the residues of $G_\rho$. If $\lambda$ has $n$ boxes,
its reverse process deletes a box of content $Y$, and
\begin{equation}\label{eq:transforms}
 G_\rho(z)=\E\left(\frac1{z-X}\,\middle|\,\rho\right),\qquad
 \frac1{G_\lambda(z)}
   =z-\alpha n\E\left(\frac1{z-(Y+d)}\,\middle|\,\lambda\right).
\end{equation}
These are the transition and cotransition formulas in
\cite[Section~2, Theorem~6.7 and Lemma~7.2]{Kerov}.
Both formulas use exactly the content and shifted inner-corner coordinates
defined above.

Both kernels have the Jack laws as coherent marginals. To check the
normalization and reversal explicitly, let $\kappa(\nu,\lambda)$ be the
positive one-box Jack Pieri multiplicity in Kerov's notation, and put
$D(\rho)=|\rho|!\alpha^{|\rho|}/c'_\rho$ and $\varphi(\rho)=1/c_\rho$.
The hook and branching formulas give
\[
 D(\lambda)=\sum_{\nu\nearrow\lambda}\kappa(\nu,\lambda)D(\nu),\qquad
 p(\nu,\lambda)=\kappa(\nu,\lambda)\frac{\varphi(\lambda)}{\varphi(\nu)},
 \qquad q(\lambda,\nu)=\kappa(\nu,\lambda)\frac{D(\nu)}{D(\lambda)}
\]
\cite[Sections~6--7]{Kerov}. Since $\pi(\rho)=D(\rho)\varphi(\rho)$,
\begin{equation}\label{eq:balance}
 \pi_{n-1,\alpha}(\nu)p(\nu,\lambda)
       =\pi_{n,\alpha}(\lambda)q(\lambda,\nu).
\end{equation}
Starting from the empty diagram and summing proves the marginal assertion
and normalization at every level.

The two transforms give the low moments needed below.

\begin{lemma}\label{lem:moments}
If $\rho$ has $m$ boxes, upward growth satisfies
\begin{equation}\label{eq:upmoments}
 \E(X\mid\rho)=0,\quad \E(X^2\mid\rho)=\alpha m,\quad
 \E(X^3\mid\rho)=2\alpha s_1(\rho)+\alpha dm.
\end{equation}
If $\lambda$ has $n$ boxes and $N=\binom n2$, reverse growth satisfies,
with $s_k=s_k(\lambda)$,
\begin{align}
 \E(Y\mid\lambda)&=\frac{2s_1}{n},\label{eq:down1}\\
 \E(Y^2\mid\lambda)&=\frac{3s_2-ds_1-\alpha N}{n},\label{eq:down2}\\
 \E(Y^3\mid\lambda)&=
 \frac{4s_3-3ds_2+d^2s_1-2\alpha(n-1)s_1+\alpha dN}{n}.
 \label{eq:down3}
\end{align}
\end{lemma}

\begin{proof}
Row-wise telescoping in \eqref{eq:G} gives
\begin{equation}\label{eq:boxproduct}
 G_\rho(z)=\frac1z\prod_{u\in\rho}
 \frac{(z-c(u))(z-c(u)-d)}{(z-c(u)-\alpha)(z-c(u)+1)}.
\end{equation}
Indeed, row $i$ of length $r$ contributes
$(z+i-1)(z+i-\alpha r)/[(z+i)(z+i-1-\alpha r)]$.
The first factors telescope over rows; the others cancel across equal
successive row lengths, leaving the genuine corners in \eqref{eq:G}.

In the following display $s_k=s_k(\rho)$.
Write $L_j$ for the coefficient of $z^{-j}$ in the formal Laurent series
$\log(zG_\rho(z))$. Its $z^{-1}$ coefficient vanishes, and a box of
content $c$ contributes
$[(c+\alpha)^j+(c-1)^j-c^j-(c+d)^j]/j$. Hence
\begin{align*}
 L_2&=\alpha m,& L_3&=2\alpha s_1+\alpha dm,\\
 L_4&=3\alpha s_2+3\alpha ds_1+(\alpha d^2+\alpha^2/2)m,\\
 L_5&=4\alpha s_3+6\alpha ds_2+(4\alpha d^2+2\alpha^2)s_1
                                      +(\alpha d^3+\alpha^2d)m.
\end{align*}
The first transform in \eqref{eq:transforms} proves \eqref{eq:upmoments}.
For the second transform take $\rho=\lambda$ and $m=n$, and expand
\[
 \frac1G=z-L_2z^{-1}-L_3z^{-2}
 -(L_4-L_2^2/2)z^{-3}-(L_5-L_2L_3)z^{-4}+O(z^{-5}).
\]
It follows that the first three moments of $Y+d$ are, respectively,
$L_3/(\alpha n)$, $(L_4-L_2^2/2)/(\alpha n)$ and
$(L_5-L_2L_3)/(\alpha n)$. Expanding $Y=(Y+d)-d$ gives
\eqref{eq:down1}--\eqref{eq:down3}. All expansions here are formal
coefficients of finite rational functions.
\end{proof}

\subsection{Three orthogonal Jack characters}

We now assume $n\ge5$, and keep $N=\binom n2$.
To estimate the cubic expression under Jack measure, use the coefficients
of power sums. Here $p_r=\sum_i x_i^r$ in formal variables and
$p_\mu=\prod_{j=1}^{\ell(\mu)}p_{\mu_j}$, with $\ell(\mu)$ the number
of parts of $\mu$. Expand the integral Jack polynomial as
$J_\lambda^{(\alpha)}=\sum_{\mu\vdash n}\theta_\mu^\lambda(\alpha)p_\mu$.
Write $\theta_k=\theta_{(k,1^{n-k})}^\lambda(\alpha)$ and
\[
 \delta=\frac{\alpha-1}{\sqrt\alpha},\quad b=\delta^2,\quad
 P_k=\frac{s_k(\lambda)}{\alpha^{k/2}}\ (k=1,2,3),\quad
 H_k=\frac{k\theta_k}{\alpha^{(k-1)/2}}\ (k=2,3,4).
\]
The character normalization makes their second moments independent of
$\alpha$. Jack coefficient orthogonality \cite[Lemmas~3.1--3.2]{Fulman}
gives
\begin{equation}\label{eq:orthogonality}
 \E H_k=0,\qquad
 \E(H_jH_k)=\mathbf1_{j=k}\,k\fall nk,\qquad
 \fall nk=n(n-1)\cdots(n-k+1).
\end{equation}
In detail, the coefficient orthogonality is
\[
 \sum_{\lambda\vdash n}\frac{\theta_\mu^\lambda\theta_\nu^\lambda}
                              {c_\lambda c'_\lambda}
    =\frac{\mathbf1_{\mu=\nu}}{z_\mu\alpha^{\ell(\mu)}},
 \qquad z_\mu=\prod_{r\ge1}r^{m_r(\mu)}m_r(\mu)!,
 \qquad \theta_{(1^n)}^\lambda=1.
\]
For $\mu=(k,1^{n-k})$, one has $z_\mu=k(n-k)!$ and
$\ell(\mu)=n-k+1$. Multiplication by $\alpha^n n!$ gives
$\E\theta_k^2=\alpha^{k-1}\fall nk/k$ and the stated means and cross terms.

Lassalle's low content identities \cite[Section~15, p.~41 of version~3]{Lassalle}
give
\begin{equation}\label{eq:Hcontents}
 \begin{aligned}
 H_2&=2P_1,\\
 H_3&=3(P_2-\delta P_1-N),\\
 H_4&=4[P_3-3\delta P_2+(2\delta^2-(2n-3))P_1+2\delta N].
 \end{aligned}
\end{equation}
His content convention is $c/\alpha$, his parameter $\beta$ is $1-\alpha$,
and his coefficient $\vartheta_k$ is $k\theta_k$. These substitutions give
\eqref{eq:Hcontents}; only the three displayed exact identities are used.
In particular $W=P_1/\sqrt N=H_2/(2\sqrt N)$ has $\E W=0$ and $\E W^2=1$.

\subsection{Conditional moments of the exchangeable pair}

Draw $\nu\sim\pi_{n-1,\alpha}$ and, conditionally on $\nu$, grow
$\lambda$ and $\lambda'$ independently by the upward kernel. Define
$W=s_1(\lambda)/\sqrt{\alpha N}$ and
$W'=s_1(\lambda')/\sqrt{\alpha N}$. By \eqref{eq:balance} these have the
required Jack marginals, and the pair is exchangeable. Conditional on
$\lambda$, its construction first removes a box of content $Y$ by the
reverse kernel, and then inserts a box of content $X$ from $\nu$.
Thus
\[
 \Delta:=W-W'=\frac{Y-X}{\sqrt{\alpha N}}.
\]

\begin{lemma}\label{lem:pair}
For this pair, $\E(\Delta\mid\lambda)=2W/n$. Moreover
\begin{align}
 A_\lambda&:=1-\frac n4\E(\Delta^2\mid\lambda)
                  =-\frac{H_3+\delta H_2}{4N},\label{eq:defect}\\
 M_\lambda&:=\E(\Delta^3\mid\lambda)
       =\frac{H_4+3\delta H_3+(5n-6+b)H_2}{nN^{3/2}}.\label{eq:cubic}
\end{align}
Their second moments are
\begin{align}
 V:=\E A_\lambda^2&=\frac{3(n-2)+2b}{4n(n-1)},\label{eq:V}\\
 \E M_\lambda^2&=\frac{8Q}{n^4(n-1)^2},\label{eq:cubicnorm}\\
 Q&=4(n-2)(n-3)+27b(n-2)+2(5n-6+b)^2.\label{eq:Q}
\end{align}
\end{lemma}

\begin{proof}
Let $D_j=\E(Y^j\mid\lambda)$, as in Lemma~\ref{lem:moments}.
Conditional on the deleted box, the inserted content has mean zero and
variance $\alpha(n-1)$.
It follows that
\[
 \E(\Delta\mid\lambda)=\frac{D_1}{\sqrt{\alpha N}}=\frac2n W,
 \qquad
 \E(\Delta^2\mid\lambda)=\frac{D_2+\alpha(n-1)}{\alpha N}
                        =\frac{3P_2-\delta P_1+N}{nN}.
\]
Formula \eqref{eq:Hcontents} gives \eqref{eq:defect}.
For the cubic moment note that $s_1(\nu)=s_1(\lambda)-Y$.
Using \eqref{eq:upmoments} conditional on the deleted box, we obtain
\begin{align*}
 \E((Y-X)^3\mid\lambda)
 &=D_3+3\alpha(n-1)D_1-2\alpha(s_1-D_1)-\alpha d(n-1)\\
 &=\frac{4s_3-3ds_2+d^2s_1}{n}+2\alpha s_1-\frac{\alpha d(n-1)}2.
\end{align*}
After division by $(\alpha N)^{3/2}$, the identity
\[
 H_4+3\delta H_3+(5n-6+b)H_2
       =4P_3-3\delta P_2+(b+2n)P_1-\delta N
\]
proves \eqref{eq:cubic}. Finally, \eqref{eq:orthogonality} gives
\[
 \E A_\lambda^2=\frac{3\fall n3+2b\fall n2}{16N^2},\qquad
 \E M_\lambda^2=
 \frac{4\fall n4+27b\fall n3+2(5n-6+b)^2\fall n2}{n^2N^3},
\]
which reduce to \eqref{eq:V} and \eqref{eq:cubicnorm}.
\end{proof}

\section{The uniform estimate}

We use the following normal-approximation result of Shao and Zhang
\cite[Corollary~2.2 of version~2]{SZ}. If $(U,U')$ is a real exchangeable
pair, $D=U-U'$, $\E(D\mid U)=\tau(U+R)$ with $0<\tau<1$, and
$\E|U|\le2$, then
\begin{equation}\label{eq:Stein}
 \dK(U,Z)\le
 \E\left|1-\frac1{2\tau}\E(D^2\mid U)\right|+\E|R|
       +2\sqrt{\frac{\E|\E(D^3\mid U)|}{\tau}}.
\end{equation}
Our pair has finite support, so all required moments exist.

\begin{proof}[Proof of Theorem~\ref{thm:main}]
First let $n\ge5$. Since $W$ is a function of $\lambda$,
Lemma~\ref{lem:pair} gives $\E(\Delta\mid W)=(2/n)W$.
Also $\E|W|\le1$, because $\E W^2=1$.
We may therefore use \eqref{eq:Stein} with $\tau=2/n$ and $R=0$.
Conditioning on $\lambda$ gives the defect $A_\lambda$ and cubic moment
$M_\lambda$ in Lemma~\ref{lem:pair}. Conditional expectation and
Cauchy--Schwarz imply
\[
 \E\left|1-\frac n4\E(\Delta^2\mid W)\right|\le\sqrt V,
 \qquad \E|\E(\Delta^3\mid W)|\le\E|M_\lambda|\le\|M_\lambda\|_2.
\]
Consequently
\begin{equation}\label{eq:reduce}
 \dK(W,Z)\le\sqrt V+2\sqrt{\|M_\lambda\|_2/(2/n)}.
\end{equation}

Put $t=\alpha/n$ and $u=b/n$. Since
$b=(\alpha-1)^2/\alpha\le\alpha$ for $\alpha\ge1$, one has $0\le u\le t$.
The exact expressions in \eqref{eq:V}--\eqref{eq:Q} give
\begin{align*}
 V&\le\frac{3+2u}{4(n-1)}\le\frac{5\max\{1,t\}}{4(n-1)},\\
 Q&\le54n^2+47bn+2b^2
    =n^2(54+47u+2u^2)\le103n^2\max\{1,t^2\}.
\end{align*}
For the last inequality, separate $u\le1$ and $u\ge1$ and bound each term
by its coefficient times $\max\{1,u^2\}$.
Substituting into \eqref{eq:cubicnorm} and \eqref{eq:reduce} yields
\[
 \dK(W,Z)\le\left(\frac{\sqrt5}{2}+2\,206^{1/4}\right)
                   \frac{\max\{1,\sqrt t\}}{\sqrt{n-1}}.
\]
Let $r=\max\{n^{-1/2},\sqrt\alpha/n\}=\max\{1,\sqrt t\}/\sqrt n$.
For $n\ge5$, $\sqrt{n/(n-1)}\le\sqrt{5/4}<9/8$.
Since $\sqrt5/2<9/8$ and $206^{1/4}<19/5$, the preceding bound is less than
\[
 \frac98\left(\frac98+\frac{38}5\right)r
       =\frac{3141}{320}r<10r.
\]
For $n=3,4$, use $\dK(W,Z)\le1$ and $r\ge n^{-1/2}\ge1/2$.
This proves \eqref{eq:main} for all the stated parameters.
\end{proof}

\paragraph{AI use.}
This manuscript was prepared with assistance from OpenAI's Codex in
literature retrieval, proof development, symbolic checks and exposition.
The mathematical claims and citations remain the author's responsibility.

\begin{thebibliography}{9}
\bibitem{CRT}
L. H. Y. Chen, M. Rai\v c and L. V. Th\`anh,
\emph{On the error bound in the normal approximation for Jack measures},
Bernoulli \textbf{27}(1) (2021), 442--468.
\href{https://doi.org/10.3150/20-BEJ1245}{doi:10.3150/20-BEJ1245}.
\bibitem{Fulman}
J. Fulman, \emph{Stein's method, Jack measure, and the Metropolis algorithm},
J. Combin. Theory Ser. A \textbf{108} (2004), 275--296.
\href{https://doi.org/10.1016/j.jcta.2004.07.003}{doi:10.1016/j.jcta.2004.07.003}.
Author version (referee revisions, June 30, 2004):
\href{https://dornsife.usc.edu/jason-fulman/wp-content/uploads/sites/224/2023/06/steinjackfinal6.pdf}{full text}.
\bibitem{Kerov}
S. Kerov, \emph{Anisotropic Young diagrams and Jack symmetric functions},
\href{https://arxiv.org/abs/math/9712267v1}{arXiv:math/9712267v1} (1997).
\bibitem{Lassalle}
M. Lassalle, \emph{Jack polynomials and free cumulants},
Adv. Math. \textbf{222} (2009), 2227--2269.
\href{https://doi.org/10.1016/j.aim.2009.07.007}{doi:10.1016/j.aim.2009.07.007}.
Author version: \href{https://arxiv.org/abs/0802.0448v3}{arXiv:0802.0448v3}.
\bibitem{SZ}
Q.-M. Shao and Z.-S. Zhang,
\emph{Berry--Esseen bounds of normal and nonnormal approximation for
unbounded exchangeable pairs}, Ann. Probab. \textbf{47}(1) (2019), 61--108.
\href{https://doi.org/10.1214/18-AOP1255}{doi:10.1214/18-AOP1255}.
Author version: \href{https://arxiv.org/abs/1710.01554v2}{arXiv:1710.01554v2}.
\bibitem{TT}
L. V. Th\`anh and N. N. Tu,
\emph{Non-uniform Berry--Esseen bounds for exchangeable pairs with
applications to the mean-field classical N-vector models and Jack measures},
C. R. Math. \textbf{363} (2025), 757--776.
\href{https://doi.org/10.5802/crmath.711}{doi:10.5802/crmath.711}.
\end{thebibliography}
\end{document}
